\chapter{Metastatic breast cancer}

\section*{Definition and Overview}
Metastatic breast cancer, also classified as Stage IV or advanced breast cancer, is defined as breast cancer that has spread beyond the primary tumour in the breast and regional lymph nodes to distant organs and tissues. The most common sites of metastasis are the bones, lungs, liver, and brain. Metastatic breast cancer can present years or decades after the treatment of early-stage disease (distant recurrence) or, less commonly, can be diagnosed at the time of first presentation (de novo metastatic disease). Although metastatic breast cancer is currently considered incurable, it is highly treatable. Recent advances in systemic therapies, including targeted agents, immunotherapies, and advanced endocrine treatments, have significantly prolonged survival, reduced symptoms, and improved quality of life, transitioning the disease into a manageable chronic condition for many patients.

\section*{Epidemiology}
Breast cancer is the most common cancer diagnosed in women worldwide and in Australia, accounting for approximately 28\% of all new female cancer diagnoses. In Australia, over 20,000 women and approximately 170 men are diagnosed with breast cancer each year. Roughly 5\% of patients present with de novo metastatic disease. Additionally, despite optimal local and systemic adjuvant therapies, approximately 20\% to 30\% of individuals diagnosed with early-stage breast cancer will eventually go on to develop metastatic recurrence. Breast cancer is the second leading cause of cancer-related mortality in Australian women, with almost all breast cancer deaths resulting from metastatic disease and subsequent organ failure.

\section*{Aetiology and Pathophysiology}
The progression to metastatic breast cancer involves a complex sequence of cellular changes known as the metastatic cascade, where cancer cells detach from the primary tumour, invade the extracellular matrix, intravasate into blood or lymphatic vessels, survive in the circulation, extravasate, and colonise distant organ parenchyma. Treatment decisions and pathophysiology are heavily dictated by the molecular and receptor subtype of the cancer.
\begin{itemize}
    \item \textbf{Hormone receptor-positive (HR+):} Cancers that express oestrogen receptors (ER) and/or progesterone receptors (PR), representing approximately 70\% of cases. These tumours are biologically less aggressive, show a strong tropism for bone tissue, and are driven by oestrogen-mediated transcription pathways, making them highly responsive to endocrine-based therapies.
    \item \textbf{HER2-positive (HER2+):} Cancers characterised by the amplification of the \textit{ERBB2} gene and overexpression of the human epidermal growth factor receptor 2 (HER2) protein, representing 15\% to 20\% of cases. Historically, these tumours carried a poor prognosis, but are now highly manageable with monoclonal antibodies and antibody-drug conjugates targeting the HER2 receptor.
    \item \textbf{Triple-negative breast cancer (TNBC):} Cancers that lack ER, PR, and HER2 expression, representing 10\% to 15\% of cases. TNBC is biologically aggressive, lacks targeted hormone options, has a high propensity for visceral (lung, liver) and central nervous system metastases, and is treated primarily with systemic chemotherapy and immunotherapy.
    \item \textbf{Germline BRCA mutations:} Inherited mutations in the \textit{BRCA1} or \textit{BRCA2} tumour suppressor genes, which impair homologous recombination DNA repair, predisposing to early-onset breast cancer and responding to specific DNA-damaging targeted therapies.
\end{itemize}

\section*{Clinical Features}
The clinical presentation of metastatic breast cancer is highly variable, depending on the specific organs affected by the metastatic burden.
\begin{itemize}
    \item \textbf{Bone metastases:} Localised, deep, aching bone pain that is typically worse at night or during weight-bearing, pathological fractures (fractures occurring from minimal trauma), and signs of spinal cord compression.
    \item \textbf{Liver metastases:} Right upper quadrant abdominal pain, abdominal distension (ascites), nausea, vomiting, loss of appetite, jaundice (yellowing of the skin and eyes), and pruritus.
    \item \textbf{Lung metastases:} Progressive dyspnoea (shortness of breath), a persistent dry cough, pleuritic chest pain, and features of pleural effusion (fluid accumulation around the lungs).
    \item \textbf{Brain metastases:} Persistent, worsening headaches (often worse in the morning or when lying down), unexplained nausea and projectile vomiting, new-onset seizures, focal neurological deficits (such as limb weakness or speech changes), and cognitive or behavioural changes.
    \item \textbf{Systemic features:} Profound fatigue, unexplained weight loss, night sweats, muscle wasting (cachexia), and low-grade fevers.
\end{itemize}

\section*{Differential Diagnosis}
In patients with a history of breast cancer presenting with new symptoms, clinicians must differentiate metastatic recurrence from other primary diseases.
\begin{itemize}
    \item \textbf{Second primary malignancy:} A new, unrelated cancer, such as a primary lung carcinoma, primary hepatocellular carcinoma, or primary osteosarcoma, which requires histological confirmation to distinguish from breast cancer metastasis.
    \item \textbf{Benign skeletal conditions:} Osteoarthritis, degenerative disc disease, osteoporotic vertebral compression fractures, or benign bone cysts, which can mimic bone metastasis pain.
    \item \textbf{Benign hepatic lesions:} Liver haemangiomas, focal nodular hyperplasia, or hepatic cysts, which can be distinguished on MRI or contrast-enhanced ultrasound.
    \item \textbf{Treatment-induced toxicities:} Radiation-induced pneumonitis, drug-induced lung injury (e.g., from interstitial lung disease associated with antibody-drug conjugates), or cardiotoxicity from anthracyclines/trastuzumab, which can mimic lung or cardiac metastasis.
\end{itemize}

\section*{When to Seek Medical Care}
Patients with a history of breast cancer should seek prompt medical advice if they notice any new, unexplained, and persistent symptoms. Acute oncological emergencies require immediate emergency department care.

\textbf{Call 000 (or local emergency services) for:}
\begin{itemize}
    \item Sudden, severe back pain accompanied by weakness or numbness in the legs, or difficulty urinating or opening the bowels (spinal cord compression).
    \item New-onset seizures or loss of consciousness.
    \item Sudden, severe shortness of breath, chest pain, or coughing up blood.
    \item High fever (above 38 degrees Celsius) accompanied by chills or shivering, particularly if the patient is currently receiving chemotherapy (neutropenic sepsis).
    \item Sudden confusion, lethargy, or extreme drowsiness.
\end{itemize}

\section*{Diagnosis and Investigations}
Confirming metastatic breast cancer requires a combination of tissue biopsy, advanced imaging, and laboratory evaluation.
\begin{itemize}
    \item \textbf{Biopsy of metastatic lesion:} Crucial to confirm the diagnosis and re-evaluate ER, PR, and HER2 status. Receptor conversion (discordance between primary and metastatic tumour receptors) occurs in up to 20\% of cases, which directly alters treatment decisions.
    \item \textbf{Staging imaging:} Computated Tomography (CT) of the chest, abdomen, and pelvis to detect visceral lesions, and a Technetium-99m bone scan to identify skeletal metastases. 18F-FDG PET-CT scans are increasingly used for highly sensitive staging.
    \item \textbf{Brain MRI:} The gold standard imaging modality for detecting and characterising brain metastases, indicated in any patient with neurological symptoms.
    \item \textbf{Laboratory evaluation:} Full blood count (to monitor for marrow infiltration or treatment cytopenias), liver function tests, renal function, serum calcium (to screen for hypercalcaemia of malignancy), and tumour markers (CA 15-3 and CEA) to monitor treatment response.
\end{itemize}

\section*{Management}
The primary objective of management is to control cancer growth, alleviate symptoms, and maintain quality of life. Systemic therapy is the mainstay of treatment, tailored to the tumour's molecular subtype.
\begin{itemize}
    \item \textbf{Hormone receptor-positive (HR+) therapy:} Managed with endocrine therapies, such as aromatase inhibitors (e.g., letrozole, anastrozole) or fulvestrant, typically combined with CDK4/6 inhibitors (e.g., ribociclib, palbociclib, or abemaciclib) as first-line therapy.
    \item \textbf{HER2-positive (HER2+) therapy:} Managed with HER2-targeted agents. First-line therapy includes trastuzumab and pertuzumab combined with a taxane chemotherapy. Second-line therapy often utilizes the highly effective antibody-drug conjugate trastuzumab deruxtecan (T-DXd).
    \item \textbf{Triple-negative breast cancer (TNBC) therapy:} Managed with systemic chemotherapy (e.g., paclitaxel, capecitabine, eribulin). Pembrolizumab (immunotherapy) is added for patients with PD-L1 positive tumours.
    \item \textbf{PARP inhibitors:} Olaparib or talazoparib are indicated for patients with germline \textit{BRCA1} or \textit{BRCA2} mutations and HER2-negative disease.
    \item \textbf{Bone-targeted therapies:} Bisphosphonates (e.g., zoledronic acid) or denosumab (RANKL inhibitor) are administered to reduce bone pain and prevent skeletal-related events, such as pathological fractures and spinal cord compression.
    \item \textbf{Local and supportive therapies:} Radiotherapy for painful bone metastases or brain metastases (stereotactic radiosurgery). Early integration of palliative care to optimise pain control, manage dyspnoea, and provide psychosocial support.
\end{itemize}

\section*{Complications}
Metastatic breast cancer can lead to life-threatening oncological complications requiring immediate medical intervention.
\begin{itemize}
    \item \textbf{Spinal cord compression:} Occurs when metastatic bone disease collapses a vertebral body, compressing the spinal cord. This is a surgical and radiotherapeutic emergency to prevent permanent paraplegia.
    \item \textbf{Pathological fractures:} Fracture of weight-bearing bones (such as the femur) due to osteolytic lesions, requiring surgical stabilisation.
    \item \textbf{Hypercalcaemia of malignancy:} Osteolysis and tumor-secreted proteins release calcium into the blood, causing confusion, renal failure, constipation, and cardiac arrhythmias.
    \item \textbf{Visceral crisis:} Severe organ dysfunction caused by high tumour burden (e.g., liver failure or respiratory failure), requiring rapid-acting chemotherapy.
    \item \textbf{Neutropenic sepsis:} Severe infection occurring in the setting of chemotherapy-induced neutropenia.
\end{itemize}

\section*{Prognosis and Outcomes}
Metastatic breast cancer is a life-limiting disease, but modern therapies have significantly extended survival times. The median overall survival is approximately three years, with a five-year survival rate of 30\% to 40\%. However, outcomes are highly variable depending on the tumour subtype. Patients with HR+ or HER2+ breast cancer can survive for five to ten years or longer due to the availability of successive lines of targeted therapies. TNBC is associated with a more aggressive disease course and a shorter median survival (typically 12 to 18 months).

\section*{Prevention and Self-Care}
Self-care strategies focus on optimizing quality of life, managing treatment side effects, and maintaining physical function.
\begin{itemize}
    \item \textbf{Treatment adherence:} Take all oral and intravenous medications precisely as prescribed, and report side effects early to allow for dosage adjustments rather than treatment discontinuation.
    \item \textbf{Physical activity:} Engage in regular, gentle exercise (such as walking or yoga) as tolerated. Exercise has been shown to reduce cancer-related fatigue, preserve muscle mass, and improve mood.
    \item \textbf{Nutritional support:} Consume a balanced, high-protein diet to prevent cachexia. Consult a dietitian if experiencing taste changes, dry mouth, or nausea.
    \item \textbf{Access support networks:} Utilise the support of breast care nurses (such as McGrath Foundation nurses), patient support groups, and psychological counseling to manage the emotional burden of living with advanced cancer.
    \item \textbf{Bone health care:} Ensure adequate calcium and Vitamin D intake, and avoid activities that carry a high risk of falls to prevent fractures.
\end{itemize}

\section*{Sources and Further Reading}
The following reputable organisations provide guidelines, patient support, and information on metastatic breast cancer:
\begin{itemize}
    \item Breast Cancer Network Australia (BCNA). \textit{Information and resources for advanced breast cancer.} https://www.bcna.org.au/
    \item Cancer Council Australia. \textit{Breast cancer clinical guidelines and patient support.} https://www.cancer.org.au/
    \item McGrath Foundation. \textit{Support from specialised breast care nurses.} https://www.mcgrathfoundation.com.au/
    \item National Comprehensive Cancer Network (NCCN). \textit{Clinical practice guidelines for breast cancer.} https://www.nccn.org/
    \item Better Health Channel. \textit{Breast cancer - metastatic.} https://www.betterhealth.vic.gov.au/health/conditionsandtreatments/breast-cancer-metastatic
\end{itemize}
